% 霍普夫代数（综述）
% license CCBYSA3
% type Wiki

本文根据 CC-BY-SA 协议转载翻译自维基百科\href{https://en.wikipedia.org/wiki/Hopf_algebra}{相关文章}

在数学中，霍普夫代数（Hopf algebra，以海因茨·霍普夫 Heinz Hopf 命名）是一种结构，它既是一个（含幺结合的）代数，又是一个（含余元的余结合）余代数；这些结构的相容性使它成为一个双代数，此外它还配备了一个满足特定性质的反同态。霍普夫代数的表示论尤其优美，因为兼容的余乘法、余元以及反元的存在，使得能够构造表示的张量积、平凡表示以及对偶表示。

霍普夫代数自然地出现在代数拓扑中（其起源之处，并与 H-空间概念相关）、群概形理论、群论（通过群环的概念），以及许多其他地方，因此可能是最常见的一类双代数。霍普夫代数也作为独立的研究对象受到关注，一方面有许多关于具体类实例的研究，另一方面也存在分类问题。它们拥有广泛的应用，从凝聚态物理与量子场论，到弦论以及大型强子对撞机（LHC）现象学。
\subsection{形式定义}
形式上，霍普夫代数是域$K$上的一个（结合且余结合的）双代数$H$，并配备一个$K$-线性映射$S: H \to H$（称为反元，antipode），使得下列图表交换：
\begin{figure}[ht]
\centering
\includegraphics[width=8cm]{./figures/080fad97a4c9a060.png}
\caption{} \label{fig_HPFds_1}
\end{figure}
这里，$\Delta$是双代数的余乘法，$\nabla$是其乘法，$\eta$ 是其单位元映射，$\varepsilon$是其余元映射。在无求和的Sweedler 记号下，这一性质也可以表述为：
$$
S(c_{(1)})c_{(2)} \;=\; c_{(1)}S(c_{(2)}) \;=\; \varepsilon(c)\,1 
\qquad \text{对于所有 } c \in H.~
$$
与代数的情况类似，在上述定义中，可以将底层的域$K$替换为交换环$R$。【4】

霍普夫代数的定义是自对偶的（这在上面图表的对称性中有所体现），因此如果可以定义 $H$的对偶（当$H$是有限维时总是可能的），那么该对偶也自动是一个霍普夫代数。【5】
\subsubsection{结构常数}
在固定了底层向量空间的一组基$\{ e_k \}$之后，可以通过乘法的结构常数来定义代数：
$$
e_i \nabla e_j = \sum_k \mu_{ij}^k e_k~
$$
对于余乘法：
$$
\Delta e_i = \sum_{j,k} \nu_i^{\,jk} \, e_j \otimes e_k~
$$
以及反元：
$$
S e_i = \sum_j \tau_i^{\,j} e_j~
$$
结合律要求：
$$
\mu_{ij}^k \, \mu_{kn}^m \;=\; \mu_{jn}^k \, \mu_{ik}^m~
$$
而余结合律要求：
$$
\nu_k^{\,ij} \, \nu_i^{\,mn} \;=\; \nu_k^{\,mi} \, \nu_i^{\,nj}~
$$
连接公理要求：
$$
\nu_k^{\,ij} \, \tau_j^{\,m} \, \mu_{im}^n 
\;=\; 
\nu_k^{\,jm} \, \tau_j^{\,i} \, \mu_{im}^n~
$$
\subsubsection{反元的性质}
反元$S$有时被要求具有一个$K$-线性逆映射，这在有限维情形下会自动成立【需要澄清】，或者当$H$是交换的、余交换的（或更一般地是拟三角的 quasitriangular）时也成立。

一般而言，$S$是一个反同态【6】，因此$S^2$是同态；若 $S$ 可逆（通常要求如此），那么 $S^2$ 就是一个自同构。

若$S^2 = \mathrm{id}_H$，则称该霍普夫代数为对合的，其底层带对合的代数是一个 $*$-代数。若$H$是有限维、在零特征域上半单、交换或余交换的，那么它是对合的。

如果一个双代数$B$承认一个反元$S$，那么这个$S$是唯一的（“一个双代数至多承认一个霍普夫代数结构”）【7】。因此，反元并不是可以随意选择的额外结构：作为霍普夫代数，是双代数的一个性质。

反元可以类比为群上的逆映射，即把元素$g$映射为$g^{-1}$【8】。
\subsubsection{霍普夫子代数}
若$H$是一个霍普夫代数，$A$是其子代数，则当 $A$ 同时也是 $H$ 的子余代数，并且反元$S$将$A$映射到$A$中时，称$A$为一个霍普夫子代数。换言之，当把 $H$的乘法、余乘法、余元和反元限制到 $A$ 上时，$A$ 本身就是一个霍普夫代数（此外还要求 $H$ 的单位元 $1$ 属于 $A$）。Warren Nichols 与 Bettina Zoeller 于 1989 年证明的 **Nichols–Zoeller 自由性定理**表明：如果 $H$ 是有限维的，则自然的 $A$-模 $H$ 是有限秩的自由模；这推广了子群的拉格朗日定理【9】。由此及积分理论推出，半单有限维霍普夫代数的霍普夫子代数自动是半单的。

在霍普夫代数 $H$ 中，若一个霍普夫子代数 $A$ 满足右稳定性条件，即

$$
\operatorname{ad}_r(h)(A) \subseteq A \quad \text{对所有 } h \in H,
$$

则称 $A$ 在 $H$ 中是**右正规**的。其中右伴随映射 $\operatorname{ad}_r$ 定义为

$$
\operatorname{ad}_r(h)(a) = S(h_{(1)})\, a \, h_{(2)}, \quad a \in A, \; h \in H.
$$

类似地，若 $A$ 在左伴随映射

$$
\operatorname{ad}_l(h)(a) = h_{(1)}\, a \, S(h_{(2)})
$$

下稳定，则称 $A$ 在 $H$ 中是**左正规**的。当反元 $S$ 是双射时，左右正规条件等价，此时称 $A$ 为一个**正规霍普夫子代数**。

在 $H$ 中的正规霍普夫子代数 $A$ 满足以下条件（作为 $H$ 的子集相等）：

$$
H A^{+} = A^{+} H,
$$

其中 $A^{+}$ 表示 $A$ 上余元的核。此正规性条件意味着 $H A^{+}$ 是 $H$ 的一个**霍普夫理想**（即：它既是余元核中的代数理想，也是余代数的余理想，并且在反元下稳定）。因此，可以得到一个商霍普夫代数

$$
H / H A^{+},
$$

以及一个满射

$$
H \to H / A^{+} H,
$$

其理论与群论中正规子群和商群的理论类似【10】。
\subsubsection{霍普夫整环}
设$R$是一个整环，$K$是其分式域。若$H$是$K$上的一个霍普夫代数，则$R$上的一个霍普夫整环$O$是$H$ 的一个整环，且在代数与余代数运算下封闭：特别地，余乘法$\Delta : O \to O \otimes O$成立【11】。
